% chapters/chapter1/sec2.tex

\section{The Irreplaceable Core Value of Integrated Sensing and Communication (ISAC)}
\label{sec:ch1-isac-core-value}

% TODO: This section is drafted by Huixin Man.
% Writing objectives:
% 1. Provide the basic definition of ISAC;
% 2. Explain the main forms of communication-sensing integration;
% 3. Briefly introduce the current development status of ISAC and its typical application scenarios.

The previous section addressed the questions of \emph{what ISAC is} and \emph{how ISAC has gradually evolved} from the perspectives of concept and technological evolution. This section proceeds to a more fundamental question: \emph{What is the core value of ISAC, and why can this value not be replaced by a simple superposition of traditional standalone systems?} We examine this question from several dimensions, including spectrum efficiency, hardware cost, deployment efficiency, and closed-loop service capability, in order to reveal the system-level advantages of ISAC over conventional standalone communication systems and standalone sensing systems. More importantly, this section argues that these advantages are not isolated local gains. Rather, they collectively point to a fundamental transformation in network capability: wireless networks are no longer limited to the single role of ``transmitting information,'' but are increasingly expected to ``sense the environment, understand system states, and support decision-making'' as part of an integrated service platform \cite{Liu2022ISAC,Gonzalez2024integrated,Saad2020Vision6G}. In this sense, ISAC should not be viewed as an optional add-on, but as one of the intrinsic capabilities of future sixth-generation (6G) networks.

\subsection{Why is ISAC not merely an ``optional addition'' to existing functions?}
\label{subsec:ch1-isac-core-value-why}
\subsubsection{From isolated functionalities to integrated system value}
\label{subsubsec:ch1-isac-core-value-logic}

Traditionally, communication and sensing systems have evolved along largely independent technological paths. Communication systems are primarily designed to deliver information reliably and efficiently, with major performance objectives including spectral efficiency, coverage, latency, and reliability~\cite{Goldsmith2005}. In contrast, sensing systems focus on understanding the physical environment, with key tasks such as target detection, parameter estimation, state tracking, and scene perception~\cite{SkolnikRadarBook}. In this sense, communication is mainly concerned with ``transmitting information,'' whereas sensing is concerned with ``observing the environment.'' For a long time, this functional separation has been both reasonable and effective.

However, many emerging 6G application scenarios, such as autonomous driving, industrial internet, low-altitude networking, and digital twins, require more than communication-only or sensing-only capabilities. Instead, they demand the tight integration of environmental awareness, information exchange, and real-time decision support within a unified system framework~\cite{Saad2020Vision6G,Gonzalez2024integrated,IMT2030Framework}. For example, in autonomous driving, a vehicle must not only sense surrounding vehicles, pedestrians, obstacles, and road conditions, but also share relevant information with other vehicles, roadside infrastructure, and edge servers through V2X links. Sensing results may directly affect communication scheduling and link configuration, while communication can in turn extend the sensing range and improve overall situational awareness~\cite{Zhang2020perceptive}. In such scenarios, communication and sensing are no longer two independent modules that can be simply combined after separate design; rather, they become mutually dependent components of a unified service loop.

If one continues to rely on a loosely coupled architecture---namely, an independent communication system, an independent sensing system, and an external decision-making module---the overall system will inevitably suffer from complex interfaces, accumulated latency, duplicated resource usage, and limited global optimization capability. Such an architecture is increasingly inadequate for future applications that require low-latency operation, high consistency, and efficient resource utilization~\cite{Cui2024integrated,Liu2020JointRadarComm}.

Therefore, the value of ISAC does not lie merely in adding a sensing function to an existing communication system, nor simply in allowing one hardware platform to perform two tasks. More fundamentally, ISAC creates a unified framework that connects sensing, communication, and service-level decision support, enabling wireless networks to evolve from pure information-delivery infrastructures into comprehensive platforms capable of environmental understanding and closed-loop task support~\cite{Liu2022ISAC,Gonzalez2024integrated}. Its significance is not that it ``does one more thing,'' but that it changes the way system capabilities are organized, thereby enabling system-level gains that are difficult to achieve with traditionally separated architectures.

\subsubsection{The Core Value of System-Level Benefit Restructuring}
\label{subsubsec:ch1-isac-core-value-benefit}

A common misunderstanding in evaluating ISAC is to judge its value by a single performance metric, such as requiring its sensing accuracy to surpass that of a dedicated radar, or its communication rate to exceed that of a standalone communication system. In practice, such a comparison is often inappropriate. A dedicated radar system is usually superior in terms of detection range, parameter estimation accuracy, target detection performance, and clutter suppression, while a dedicated communication system may achieve better throughput, error-rate performance, and link optimization capability \cite{Goldsmith2005,Liu2022ISAC,SkolnikRadarBook}. This is because a standalone system can devote all available resources to one function, whereas an ISAC system must support both communication and sensing under a unified resource constraint. Its design is therefore essentially a problem of multi-objective tradeoff and joint optimization.

Hence, the core value of ISAC does not primarily lie in outperforming dedicated systems in every individual function, but in restructuring the overall system-level benefit. Rather than separately pursuing ``communication optimality'' or ``sensing optimality,'' ISAC seeks to achieve \emph{overall service effectiveness} on a unified platform \cite{Liu2022ISAC,Gonzalez2024integrated}. Its system-level benefits can be summarized in the following aspects.

\begin{enumerate}[label=(\arabic*)]
\item \textbf{Improved resource utilization.}  
Communication and sensing can share spectrum, hardware platforms, antenna arrays, and time--frequency resources. This reduces resource redundancy, alleviates spectrum scarcity, and improves overall utilization efficiency at the system level \cite{Liu2020JointRadarComm,Liu2022ISAC,Hassanien2015dual}.

\item \textbf{Reduced total system cost.}  
Through hardware sharing, site reuse, and infrastructure integration, ISAC can lower the overall cost of deployment and operation. Although the design of an integrated system may be more complex, it is often more advantageous from the perspective of total cost of ownership and large-scale deployment \cite{Liu2022ISAC,Wen2024survey}.

\item \textbf{Greater deployment flexibility and scalability.}  
ISAC can build upon existing communication infrastructures, such as cellular networks, Wi-Fi systems, and vehicular networks, and introduce sensing capability through protocol upgrades, waveform redesign, and algorithm enhancement. This enables rapid service rollout and supports the evolutionary expansion of network functionality \cite{Ma2022WiFiSensingSurvey}.

\item \textbf{A shift from connectivity to closed-loop intelligence.}  
ISAC enables a unified ``sensing--communication--decision--feedback'' loop. In this framework, the network not only delivers information but also supports environmental awareness, cooperative control, and intelligent decision-making. Such a capability is particularly important in applications such as autonomous driving, industrial automation, multi-robot collaboration, and digital twins \cite{Gonzalez2024integrated,Saad2020Vision6G}.
\end{enumerate}

Therefore, the value of ISAC should not be assessed solely by whether a single metric exceeds that of a dedicated system. A more appropriate evaluation should consider its overall system-level benefits in terms of resource efficiency, total cost, deployment evolution, and support for closed-loop services. The key significance of ISAC lies not in maximizing communication and sensing separately, but in redefining what constitutes a better system on a unified platform.

\subsection{Spectrum Efficiency Advantage}
\label{subsec:ch1-isac-core-value-spectrum}
\subsubsection{Spectrum Occupation in Traditional Separated Systems}
\label{subsubsec:ch1-isac-core-value-spectrum-occupation}

Spectrum is the fundamental physical medium of wireless systems and one of the scarcest strategic resources in modern information society. Communication, radar, navigation, remote sensing, and industrial wireless control all rely on access to limited electromagnetic spectrum resources to perform their core functions \cite{Goldsmith2005}. Traditionally, communication and sensing systems have been assigned separate spectrum bands, and mutual interference has been avoided mainly through static or semi-static frequency partitioning. Although this approach simplifies regulation, hardware design, and interference management, it becomes increasingly inefficient as wireless services grow in density and heterogeneity \cite{Gonzalez2024integrated}.

In practice, radar systems have long occupied several important bands, such as the L, S, C, and X bands, while the demand for spectrum from cellular networks, wireless local area networks, satellite Internet, and emerging wireless access services continues to increase \cite{SkolnikRadarBook}. Under the evolution toward 5G-Advanced and 6G, wireless networks are expected to support not only higher data rates, lower latency, and massive connectivity, but also environmental sensing capabilities \cite{Liu2022ISAC,Saad2020Vision6G,IMT2030Framework}. As a result, the traditional paradigm of exclusive and isolated spectrum allocation is becoming less suitable for future wireless systems.

More importantly, spectrum assignment does not necessarily imply efficient spectrum utilization. Many radar systems, especially pulsed radars, exhibit intermittent transmission, directional beamforming, and task-dependent operation, which means that their spectrum occupancy is often discontinuous in time, space, and mission state \cite{SkolnikRadarBook}. Communication systems also show strong spatiotemporal traffic imbalance: spectrum usage can vary significantly across regions and time periods, leaving some allocated resources underutilized in specific scenarios \cite{Haykin2005CognitiveRadio}. Therefore, the main challenge is not only the scarcity of total spectrum, but also the lack of cross-system coordination, which leads to fragmentation of spectrum usage across time, frequency, and space \cite{Liu2020JointRadarComm,Liu2022ISAC}.

This problem becomes even more pronounced in the millimeter-wave and terahertz regimes. As both communication and sensing systems expand into higher-frequency wideband spectrum, their demands increasingly overlap in bands such as 24~GHz, 28~GHz, 39~GHz, and 60--77~GHz \cite{Liu2022ISAC,Gonzalez2024integrated}. This overlap is particularly evident in vehicular and autonomous driving scenarios, where communication links and automotive radars often operate in nearby frequency ranges and are tightly coupled in terms of deployment location, coverage area, and service targets \cite{Sturm2011WaveformDesign,Kumari2018IEEE80211adRadar}. In such cases, rigid frequency separation not only reduces spectrum efficiency but also limits the potential for deeper system cooperation.

From the ISAC perspective, the key is therefore not simply to acquire more dedicated spectrum, but to move beyond the conventional paradigm of exclusive occupation toward one of on-demand sharing, joint design, and dynamic reuse \cite{Liu2020JointRadarComm,Liu2022ISAC,Gonzalez2024integrated}. Only by placing communication and sensing within a unified time-frequency-space resource framework can spectrum be transformed from a statically partitioned scarce asset into a dynamically orchestrated integrated capability. This forms the starting point for understanding the spectrum efficiency advantage of ISAC.

\subsubsection{Sources of ISAC Advantages in Spectrum Sharing and Reuse}
\label{subsubsec:ch1-adv source of isac-spectrum}

The spectrum-efficiency gain of ISAC does not simply come from forcing communication and sensing to operate in the same frequency band at the same time. Without unified design, such naive coexistence may even lead to severe mutual interference and performance degradation. The key advantage of ISAC lies in the joint organization and reuse of multidimensional resources, including time, frequency, space, and code, within a unified system framework, so that the same spectrum can simultaneously support both information transmission and environmental sensing. In general, this advantage comes from three main aspects: functional reuse of the same resources, joint signal design for dual purposes, and network-level dynamic spectrum coordination.

\paragraph{1) Functional reuse of the same time-frequency resources.}
The most direct spectrum gain comes from using the same time-frequency resources for both communication and sensing. In conventional systems, communication and radar usually occupy separate frequency bands and transmission slots. In contrast, in an ISAC system, the same transmitted signal can simultaneously carry communication data and serve as a sensing waveform \cite{Liu2022ISAC,Hassanien2015dual,Sturm2011WaveformDesign}. For example, in OFDM-based ISAC, one OFDM symbol can convey modulated information to a communication receiver, while its echoes reflected from targets can be used for estimating range, Doppler, or angle. In multi-antenna systems, the same transmit beam can provide directional communication gain and, at the same time, illuminate a target region for sensing \cite{Liu2020JointRadarComm,Gonzalez2024integrated}. As a result, resource blocks that would otherwise be separately allocated to communication and sensing can now support both functions simultaneously, reducing the overall spectrum requirement.

A simple example helps illustrate this point. Suppose that a communication system requires about 100~MHz bandwidth to support a target data rate, while a sensing system also requires about 100~MHz effective bandwidth to achieve a desired range resolution. Under conventional separate deployment, the total spectrum demand may be around 200~MHz. Under an ideal ISAC sharing mode, however, the same 100~MHz signal may support both tasks. This example is only intended to show the directional benefit of sharing; it does not imply that practical systems can always reduce bandwidth demand by half without loss. In real deployments, communication and sensing often have different preferences in waveform structure, power allocation, pilot usage, and performance metrics \cite{Liu2020JointRadarComm,Liu2022ISAC}. Even so, functional reuse remains one of the most fundamental sources of ISAC spectrum gains.

\paragraph{2) Deeper spectrum gains enabled by joint waveform and signal design.}
If functional reuse means that the same resource can perform two tasks, then joint signal design determines how the same resource can perform both tasks more efficiently. A deeper spectrum advantage of ISAC comes from the joint optimization of waveform, pilot structure, frame format, beamforming, and power allocation. Instead of simply superimposing two independently designed signals, ISAC considers communication and sensing requirements from the beginning of signal design.

OFDM is a representative example because its multicarrier structure is naturally compatible with both communication and sensing. On the communication side, OFDM supports flexible subcarrier modulation, equalization, and multiuser access. On the sensing side, its wideband frequency-domain structure is useful for range estimation, while phase evolution across OFDM symbols can support velocity sensing. By properly designing pilot patterns, subcarrier mapping, cyclic prefix length, and frame structure, reference signals originally introduced for synchronization or channel estimation in communication can also be reused for sensing, allowing target information to be extracted without requiring a fully separate sensing spectrum allocation \cite{Ma2022WiFiSensingSurvey,Kumari2018IEEE80211adRadar}. Similarly, in MIMO-ISAC systems, spatial beams not only serve communication coverage and interference control, but also determine sensing illumination regions and angular resolution, thereby improving the joint utilization of spectral and spatial resources \cite{Liu2020JointRadarComm,Liu2022ISAC}.

\paragraph{3) Dynamic spectrum coordination at the network level.}
Beyond resource reuse and joint signal design on a single link, the spectrum advantage of ISAC also extends to network-level scheduling. In future cellular networks, vehicular networks, industrial wireless systems, and low-altitude collaborative networks, both communication traffic and sensing demand are highly time-varying and location-dependent \cite{Saad2020Vision6G,Cui2024integrated,IMT2030Framework}. Since communication and sensing share a common platform in ISAC, the system can jointly schedule time slots, frequency bands, beams, transmit power, and even computational resources according to scenario conditions and service objectives \cite{Zhang2020perceptive,Cui2024integrated}. For instance, when communication traffic is light, more beam dwelling time or denser reference signaling can be assigned to sensing; when sensing demand is weak, more resources can be prioritized for data transmission. Compared with conventional separate systems, such dynamic coordination transforms fragmented and statically partitioned spectrum into a unified resource pool that can be flexibly allocated according to task needs.

In summary, the spectrum-efficiency advantage of ISAC does not come merely from frequency sharing itself. Rather, it arises from multi-layer coordination: dual use of the same resources at the resource level, joint waveform and signal optimization at the signal level, and dynamic cross-task scheduling at the network level. These mechanisms together explain why ISAC can achieve spectrum sharing and reuse gains that cannot be obtained by simply combining conventional communication and sensing systems.

\subsubsection{Applicable Conditions and Boundaries of the Spectral Efficiency Advantage}
\label{subsubsec:ch1-isac-core value-spectrum-boundary}

Although ISAC has the potential to improve spectral efficiency, this advantage does not arise automatically whenever communication and sensing share the same spectrum. Rather, it is a conditional gain that depends on the application scenario, system architecture, and the quality of joint design. When the requirements of communication and sensing are strongly conflicting, or when transceiver capability, interference suppression, and resource scheduling are inadequate, the theoretical benefit of spectrum sharing may not translate into practical performance gains. In some cases, the cost of sharing may even outweigh its benefit.

\paragraph{1) The inherent tension between communication-oriented and sensing-oriented signal requirements.}
Communication and sensing do not impose identical requirements on signal design. Communication systems generally favor signals that can carry as much information as possible in order to maximize spectral efficiency, throughput, and multi-user multiplexing capability. By contrast, sensing systems often prefer waveforms with desirable autocorrelation, cross-correlation, and ambiguity-function properties to support reliable detection, estimation, and tracking \cite{SkolnikRadarBook}. Therefore, when a single signal is expected to serve both purposes, the system must balance multiple objectives, including rate, error performance, detection probability, and estimation accuracy. This tension becomes especially pronounced in extreme cases, such as heavily loaded communication links or high-resolution sensing tasks, where the benefit of spectrum sharing may be significantly reduced.

\paragraph{2) Self-interference and transceiver implementation constraints.}
In monostatic ISAC systems, simultaneous transmission and reception over the same frequency band leads to severe self-interference, similar to the challenge encountered in full-duplex communications \cite{Sabharwal2014InBandFullDuplex}. Since transmitter leakage is often much stronger than the echo from distant or weak targets, insufficient analog isolation, digital cancellation, or hardware impairment compensation can severely compress the dynamic range of the sensing receiver. As a result, weak-target detection and parameter estimation may degrade substantially \cite{Gonzalez2024integrated,Bliss2007MIMOFullDuplex}. Hence, the spectral efficiency gain of ISAC often comes at the price of higher transceiver complexity and more stringent hardware requirements.

\paragraph{3) Sensing requirements impose additional constraints on resource allocation.}
Sensing performance is jointly determined by several physical factors, including bandwidth, observation time, transmit power, and array aperture. In general, finer range resolution requires larger bandwidth, finer velocity resolution requires longer coherent observation time, and finer angular resolution requires a larger aperture or more refined beam control \cite{Richards2014,SkolnikRadarBook}. However, these requirements do not always align with the resource allocation strategies preferred in communication systems, which often emphasize low latency, flexible scheduling, and efficient service to multiple users. Therefore, the spectral efficiency advantage of ISAC must always be evaluated under explicit sensing-performance constraints rather than in isolation from the sensing task itself.

\paragraph{4) Interference coupling becomes more complicated in multi-user and multi-target scenarios.}
In multi-user, multi-target, and multi-base-station scenarios, the benefit of spectrum sharing is further limited by more complex interference coupling \cite{Liu2022ISAC,Cui2024integrated}. Communication links suffer from multi-user interference, while sensing tasks are affected by clutter, sidelobe leakage, and parameter coupling. These effects also interact through shared beamforming, power allocation, and spectrum occupancy. As system scale increases, simple spectrum reuse strategies often become insufficient, and more sophisticated mechanisms are required, such as joint beamforming, power control, pilot design, and task scheduling \cite{Liu2020JointRadarComm,Gonzalez2024integrated}. Therefore, the spectral efficiency advantage of ISAC does not necessarily grow linearly with system size and may instead be constrained by increasingly difficult interference management.

In summary, the spectral efficiency advantage of ISAC should be understood as a conditional system-level benefit rather than a universal conclusion. It can be fully realized only when communication and sensing requirements are sufficiently compatible, self-interference can be effectively suppressed, sensing-resolution demands can be coordinated with communication scheduling, and multi-user/multi-target interference can be properly managed. For this reason, later discussions on waveform design, transceiver design, and resource optimization should be based on quantitative modeling and boundary analysis, rather than on purely conceptual claims.

\subsection{Hardware Cost Advantage}
\label{subsec:ch1-isac-core value-hardware}
\subsubsection{The Problem of Hardware Duplication in Standalone Deployment}
\label{subsubsec:ch1-isac-core value-hardware-duplication}

When communication and sensing are carried out by two independent systems, it is usually necessary to build two relatively complete hardware chains in parallel, including antenna arrays, RF front-ends, analog-to-digital/digital-to-analog converters (ADC/DAC), baseband processing units, clock and local-oscillator modules, power supply and thermal management units, as well as installation platforms such as cabinets, towers, or vehicular mounting structures. In large-scale network deployment scenarios, such duplicated construction significantly increases both capital expenditure (CAPEX) and operating expenditure (OPEX) \cite{Zhang2021ISAC6G,Liu2022ISAC,ITU2030}.

From the perspective of system composition, traditional communication equipment and traditional sensing equipment differ in their functional objectives, but they share a substantial number of similar hardware modules. For example, both cellular base stations and mmWave radar systems rely on high-frequency RF front-ends, antenna arrays, beamforming networks, and high-speed digital signal processors. Under a standalone deployment scheme, these high-value components are often purchased, installed, powered, and maintained separately, while also occupying separate sites, spectrum, backhaul resources, and maintenance effort.

Consider the V2X scenario as an example. If each roadside node deploys both a communication device (e.g., a cellular V2X (C-V2X) roadside unit) and a dedicated sensing device (e.g., mmWave radar, camera, or light detection and ranging (LiDAR)), then at least the following forms of duplicated investment will arise:
\begin{itemize}
    \item \textbf{Duplicated front-end hardware:} communication and sensing devices each require their own antennas, RF chains, clock sources, and processors;
    \item \textbf{Duplicated infrastructure:} each system requires independent installation space, power interfaces, thermal design, and protective enclosures;
    \item \textbf{Duplicated network access:} each system requires its own data backhaul and network management interface;
    \item \textbf{Duplicated operation and maintenance:} each system requires separate monitoring, calibration, maintenance, and fault diagnosis processes.
\end{itemize}

In a medium-sized city, the number of nodes at intersections, roadside poles, parking areas, and tunnel entrances can easily reach several thousand or more. In such a case, the cost difference at a single node is amplified by network scale, making the question of whether hardware platforms are shared no longer merely a technical issue, but one that directly determines the economic sustainability of system deployment.

Beyond the initial construction cost, standalone deployment also significantly increases lifecycle cost. First, two sets of equipment typically imply higher total power consumption. Second, a larger number of devices increases the probability of faults and complicates spare-parts management. Third, sensing devices often require periodic calibration and environmental adaptation, which further raises maintenance costs. Therefore, in large-scale scenarios such as 6G, intelligent transportation, smart campuses, and the industrial Internet, reducing hardware duplication has become one of the practical drivers behind the development of ISAC.

\subsubsection{The Basic Logic of Hardware Sharing in ISAC}
\label{subsubsec:ch1-isac-core value-hardware-sharing}

The hardware cost advantage of ISAC fundamentally originates from the idea that \textbf{the same set of physical resources serves both communication and sensing simultaneously}. Such sharing may occur at multiple levels, including front-end hardware, platform infrastructure, signal-processing chains, and network control.

\paragraph{1) Front-end sharing}
At the front-end level, communication and sensing can share antenna arrays, beamforming networks, RF transceiver chains, and part of the clock/local-oscillator resources. This is particularly natural at mmWave and THz frequencies, where both communication and sensing rely heavily on narrow-beam gain to overcome severe path loss. As a result, they have inherently common requirements for large-scale arrays, hybrid beamforming architectures, and wideband RF front-ends \cite{Zhang2021ISAC6G, Wei2023Integrated}.  
For example, in a monostatic ISAC base station, the same phased array can generate communication beams toward users and sensing illumination beams toward the environment; the same power amplifiers, low-noise amplifiers, and mixer chains can also support dual-function operation under unified scheduling.

\paragraph{2) Baseband and algorithmic reuse}
At the baseband level, communication and sensing exhibit substantial similarity in algorithmic structure. Taking OFDM as an example, modules such as synchronization, fast Fourier transform/inverse fast Fourier transform (FFT/IFFT), channel estimation, pilot processing, and beam training in communication systems are mathematically related to delay estimation, Doppler estimation, angle estimation, and environmental representation in sensing systems \cite{Liu2022ISAC,Sturm2011WaveformDesign}.  
Therefore, ISAC does not necessarily require a completely separate digital processing platform for sensing. In many cases, sensing capability can be added to the existing communication baseband processor through software upgrades, task scheduling, and limited hardware enhancement.

\paragraph{3) Platform-level integration}
At the platform level, existing cellular base stations, Wi-Fi access points, roadside units (RSUs), UAV platforms, and in-vehicle terminals can all serve as natural carriers for ISAC. Consider a cellular base station as an example. It already possesses several key features: multi-antenna or massive MIMO arrays; wideband RF and high-speed baseband processing capability; stable power supply, thermal conditions, and installation support; backhaul links to the edge cloud or core network; and built-in network management and clock synchronization mechanisms.  
These capabilities overlap strongly with the hardware foundations required for high-resolution environmental sensing. Therefore, compared with constructing dedicated sensing infrastructure from scratch, introducing sensing capability on top of existing communication platforms usually incurs a much lower marginal cost.

\paragraph{4) Network and data sharing}
The cost advantage of ISAC is not limited to the sharing of transmission and reception equipment; it also extends to the sharing of information links. Traditional standalone sensing systems often require dedicated data backhaul, target fusion, and control interfaces. In ISAC, by contrast, sensing results can be uploaded directly through the communication network to edge servers or control centers, while the communication links themselves can carry sensor data, target states, and collaborative control information. This reduces the need for additional networking infrastructure.

\paragraph{5) A useful cost decomposition}
For pedagogical clarity, the hardware cost advantage of ISAC can be summarized in the forms listed in Table~\ref{tab:isac-hardware-saving}.

\begin{longtable}{p{3.3cm}p{4.8cm}p{5.2cm}}
\caption{Main Sources of Hardware Cost Saving in ISAC}
\label{tab:isac-hardware-saving} \\
\hline
\textbf{Layer} & \textbf{Shared Resources} & \textbf{Cost Benefit} \\
\hline
\endfirsthead

\multicolumn{3}{c}{\tablename~\thetable{} -- continued from previous page} \\
\hline
\textbf{Layer} & \textbf{Shared Resources} & \textbf{Cost Benefit} \\
\hline
\endhead

\hline
\multicolumn{3}{r}{Continued on next page} \\
\endfoot

\hline \hline
\endlastfoot

RF front-end & Antennas, RF chains, oscillators, beamforming modules & Reduce duplicated high-value hardware \\
Baseband processing & FFT/IFFT, synchronization, channel/parameter estimation modules & Lower processor and accelerator redundancy \\
Infrastructure platform & Towers, poles, cabinets, power supply, cooling & Save site acquisition and installation cost \\
Networking & Backhaul, edge access, control interface & Reduce networking and integration overhead \\
Operations \& maintenance & Monitoring, fault diagnosis, software upgrade & Lower lifecycle OPEX \\
\end{longtable}

It should be emphasized that ``sharing'' in ISAC does not mean that all modules must be completely unified. A more common engineering approach is to \textbf{share the main platform and the expensive common components, while adding a limited number of dedicated modules for sensing}, such as more stable clocks, larger buffers, dedicated calibration loops, or higher-resolution data-acquisition units. Even so, as long as the incremental cost remains lower than that of building a complete standalone sensing system, the overall solution is still economically attractive \cite{Liu2022ISAC, Wei2023Integrated}.

\subsubsection{Limitations of Hardware Sharing and the Associated Engineering Costs}
\label{subsubsec:ch1-isac-core value-hardware-limitations}

Although hardware sharing can significantly reduce overall investment, such sharing is not free of cost. The fundamental reason is that communication and sensing do not have perfectly aligned performance objectives, and in some hardware dimensions they even exhibit clear tension. Therefore, the cost advantage of ISAC should be understood as a \textbf{reduction in total system cost}, rather than as ``every individual component becoming cheaper.''

\begin{itemize}
\item \textbf{Higher dynamic-range requirements:}  
In radar-like sensing, the target echo power is usually much lower than the transmit signal, and the echo strengths corresponding to different targets, distances, and radar cross sections may vary greatly. The receiver must maintain stable operation in the presence of strong leakage, strong short-range echoes, and weak long-range echoes, which often requires a higher dynamic range. In contrast, although communication receivers also need to handle received-power variations across users, they are not typically designed around extremely weak echo detection. If an ISAC platform must satisfy both requirements, then ADC resolution, automatic gain control, and front-end linearity may all need enhancement, leading to additional cost \cite{Wei2023Integrated}.

\item \textbf{More difficult transmit--receive isolation and self-interference suppression:}  
In monostatic ISAC, especially under simultaneous transmission and reception, transmit leakage, self-interference, and platform clutter may severely mask target echoes. To ensure reliable sensing, strong antenna isolation, analog-domain cancellation, and digital-domain interference suppression are often required. These measures increase the complexity of antenna layout, RF chains, and calibration circuits \cite{Sturm2011WaveformDesign, Hassanien2016signaling}.

\item \textbf{Stricter clock stability and phase-noise constraints:}  
High-accuracy range, Doppler, and micro-Doppler estimation are typically more sensitive to carrier stability, sampling-clock jitter, and phase noise. Although communication systems are also affected by phase noise, in many cases the effects can be mitigated through synchronization, pilots, and equalization. In sensing, however, fine velocity estimation, human-activity recognition, and high-resolution parameter estimation often rely more critically on stable local oscillators. Therefore, shared clock design in ISAC usually needs to satisfy the stricter precision requirement of the sensing side \cite{Sen2013YouAreWhat}.

\item \textbf{Tradeoff between linearity and efficiency in power amplifiers:}  
Communication systems, especially those using high-peak-to-average-power-ratio (high-PAPR) waveforms such as OFDM, require high linearity from power amplifiers. Traditional radar, on the other hand, often emphasizes high efficiency and high-energy transmission, and may prefer constant-modulus or low-PAPR waveforms. ISAC waveforms and hardware must jointly account for communication-friendliness and sensing-friendliness, making the coordinated optimization of power amplifiers, digital predistortion, and waveform design even more important \cite{Sturm2011WaveformDesign, Hassanien2016signaling}.

\item \textbf{Calibration and consistency-maintenance cost cannot be ignored:}  
Shared arrays and shared RF chains reduce the number of hardware modules, but they also make the system more dependent on precise calibration. Amplitude/phase errors across array elements, channel inconsistency, temperature drift, and mechanical misalignment affect not only communication beamforming but also sensing angle estimation and imaging quality. Therefore, large-scale ISAC deployment often requires more systematic online calibration and health-monitoring mechanisms, which translate into additional engineering complexity and software maintenance cost.

\item \textbf{Not all sensors can be fully replaced:}  
It should be pointed out that the hardware-sharing philosophy of ISAC mainly applies to \textbf{radio-based sensing}, such as range, velocity, angle, occupancy detection, and coarse-grained environmental awareness. For tasks that require very rich semantic detail, texture, color, or highly accurate three-dimensional (3D) contour information, cameras and LiDAR still possess irreplaceable advantages. Therefore, in autonomous driving, robotics, or industrial inspection, ISAC is more likely to become a key component within a multisensor system, rather than completely replacing all other sensing hardware in every scenario \cite{Wei2023Integrated}.
\end{itemize}

In summary, hardware sharing in ISAC is not ``zero-cost sharing.'' Rather, it trades \textbf{a moderate increase in the design complexity of a unified platform} for \textbf{network-level, system-level, and lifecycle-level cost reduction}. From the perspective of engineering economics, as long as the CAPEX/OPEX savings brought by sharing exceed the additional costs of devices, calibration, and algorithms required to support dual-function operation, ISAC has a clear hardware-economic advantage. Later chapters will further discuss hardware nonidealities, transceiver architectures, and practical engineering implementation issues for real deployment.

\subsection{Deployment Efficiency Advantage}
\label{subsec:ch1-isac-core value-deployment}
\subsubsection{Why Does the Future Network Require High Deployment Efficiency}
\label{subsubsec:ch1-isac-core value-deployment-why}

The objective of future networks is no longer limited to providing high-speed connectivity. Instead, networks are evolving toward an integrated information infrastructure that jointly supports communications, sensing, computing, and control. For emerging scenarios such as 6G, intelligent transportation, the low-altitude economy, industrial Internet, and smart cities, future networks must not only deliver higher performance, but also enable faster deployment, greater scalability, and more flexible upgrading.

On the one hand, future application scenarios are becoming broader, more dynamic, and more diverse in their requirements. For example, low-altitude applications require networks to support both wireless connectivity for unmanned aerial vehicles and real-time sensing of airspace targets and risks. In vehicle-to-everything and road-side cooperation scenarios, infrastructure is expected to provide both information exchange and environmental sensing. In industrial Internet scenarios, networks must simultaneously support highly reliable connectivity, low-latency coordination, and real-time perception of production environments \cite{Zhang2021ISAC6G, Feng2020JointRadarCommunication, Shafi20175G}. Since these scenarios are often characterized by wide coverage, numerous sites, complex environments, and rapidly changing service demands, long deployment cycles can significantly hinder service rollout and large-scale adoption.

On the other hand, the traditional mode in which communication systems and sensing systems are deployed separately is not well suited to the agile deployment required by future networks. Under this approach, communication infrastructure and sensing infrastructure must be planned, installed, connected, and maintained independently. This often leads to duplicated investment in sites, power supply, backhaul, and maintenance, as well as increased system heterogeneity and longer construction cycles. Such limitations become even more evident in highly dynamic scenarios, such as temporary event support, post-disaster emergency communications, and short-term traffic management.

Meanwhile, network evolution is moving toward software-defined, virtualized, and platform-based architectures. Since 5G, an increasing number of network functions can be introduced through software upgrades, cloud deployment, and open interfaces rather than through extensive new hardware installation \cite{Andrews2014, Letaief2019Roadmap}. In this context, ISAC provides a deployment-friendly solution by embedding sensing capability into existing communication infrastructure. If sensing functions can be activated through software upgrades, parameter reconfiguration, and limited hardware enhancement, new capabilities can be introduced without substantially increasing deployment complexity.

Therefore, future networks require high deployment efficiency because application expansion is rapid, service demands change frequently, and the traditional separately deployed architecture is both costly and time-consuming. By integrating sensing into communication infrastructure, ISAC offers an important technical path toward faster deployment, greater service agility, and improved operational efficiency \cite{Liu2022ISAC, IMT2030Framework,ITU2030}.

\subsubsection{Mechanisms Through Which ISAC Improves Deployment Efficiency}
\label{subsubsec:ch1-isac-core value-deployment-mechanism}

The core mechanism through which ISAC improves deployment efficiency can be summarized as follows: \textbf{leveraging existing communication infrastructure, network capabilities, and operation-and-maintenance systems to rapidly add sensing functionality at low incremental cost}. This mechanism not only reduces duplicated hardware construction, but also shortens the entire chain of site planning, installation and commissioning, network access, and service activation.

\paragraph{1) Reuse of existing communication infrastructure}
Current cellular networks have already formed a mature infrastructure with broad coverage and dense site distribution. Many of the capabilities provided by this infrastructure overlap strongly with the basic requirements of wireless sensing. Therefore, if a base station can be enabled---through software upgrade, waveform reconfiguration, and limited hardware expansion---to provide environmental sensing while continuing to offer communication service, then the construction of a sensing network no longer needs to start from the full sequence of ``site application--equipment procurement--installation--integration and debugging.'' Instead, sensing can be quickly overlaid on the existing network.

From an engineering perspective, this approach brings at least four deployment benefits:
\begin{itemize}
    \item \textbf{Reduced need for new sites:} sensing functionality can be attached to existing sites, avoiding the need to build large numbers of independent sensing locations;
    \item \textbf{Reduced civil-engineering and support cycles:} existing power supply, cooling, cabinets, towers, and backhaul resources can be directly reused;
    \item \textbf{Reduced system-integration workload:} sensing functions can be connected to existing network management, orchestration, and edge platforms;
    \item \textbf{Reduced approval and coordination complexity:} fewer new devices mean fewer additional site approvals, construction coordination steps, and operation-and-maintenance access procedures.
\end{itemize}

\paragraph{2) Unified planning for newly deployed nodes}
For newly built networks or newly expanded areas, ISAC can also improve deployment efficiency through ``unified planning, unified construction, and unified activation.'' If communication and sensing are carried on the same platform, then site selection, coverage planning, power design, backhaul access, edge processing, and network management can all be completed within the same engineering workflow, rather than being separately planned and built for communication and sensing.  
The advantage of such unified planning lies not only in shortening the construction cycle, but also in reducing cross-system coordination and the difficulty of later retrofitting. This is especially beneficial in scenarios such as roads, campuses, ports, airports, and mines, where unified deployment is often more feasible than ``building communication first and retrofitting sensing later.''

\paragraph{3) Faster service rollout}
From the perspective of service activation, ISAC significantly lowers the threshold for introducing new sensing services. Under the traditional model, if a certain area requires new capabilities such as UAV supervision, traffic-state monitoring, or perimeter environmental awareness, it is usually necessary to re-plan dedicated sensing nodes, deploy backhaul links, and establish separate data processing and operation workflows. Under the ISAC model, however, operators can first activate relevant functions on existing communication nodes, and then bring the capability online through software updates, parameter configuration, and limited hardware enhancement. In this sense, ISAC makes the addition of network capability more like a ``platform upgrade'' than a ``network rebuild.'' This is particularly important for emerging services that require rapid pilot deployment, rapid replication, and rapid scaling.

\paragraph{4) Deployment efficiency enabled by software-defined evolution}
Another important deployment advantage of ISAC is its strong consistency with the software-defined evolution path of future networks. As wireless networks increasingly rely on software configuration, cloud-native baseband processing, open interfaces, and intelligent orchestration, many functions that previously required new dedicated devices are being transformed into functions achieved through software reconfiguration and resource scheduling \cite{Andrews2014,Letaief2019Roadmap}.  
Under this trend, if sensing can be provided as an intrinsic network function, then its activation process becomes more similar to the extension of communication functions, namely: on-demand activation, region-based enablement, policy configuration, and progressive upgrade. This makes ISAC not only superior to traditional standalone sensing solutions in deployment efficiency, but also more aligned with the future direction of network cloudification and hardware--software decoupling.

For teaching purposes, the major mechanisms through which ISAC improves deployment efficiency can be summarized as shown in Table~\ref{tab:isac-deployment-mechanism}.

\begin{table}[htbp]
\centering
\caption{Mechanisms Through Which ISAC Improves Deployment Efficiency}
\label{tab:isac-deployment-mechanism}
\begin{tabular}{p{3.5cm}p{5.0cm}p{4.8cm}}
\hline
\textbf{Mechanism} & \textbf{Key Idea} & \textbf{Deployment Benefit} \\
\hline
Infrastructure reuse & Reuse sites, antennas, power supply, backhaul and cabinets & Shorter construction cycle \\
Unified planning & Joint planning of communication and sensing functions & Lower coordination complexity \\
Software upgrade & Enable sensing through software/configuration updates & Faster service rollout \\
Platform integration & Use existing network management and edge platforms & Easier operation and scaling \\
Incremental enhancement & Add only limited hardware when needed & Lower deployment threshold \\
\hline
\end{tabular}
\end{table}

Of course, the above advantages usually depend on one important premise: the communication platform itself must possess a certain sensing potential, such as sufficient bandwidth, adequate array size, clock stability, and computing resources. If the target application places extremely high requirements on sensing accuracy, detection range, or target-type recognition capability, then dedicated modules or even other sensing devices may still be required. Therefore, the deployment-efficiency advantage of ISAC is essentially a \textbf{system-level advantage of extending sensing capability as rapidly as possible by maximally leveraging existing network platforms}, rather than a guarantee that no additional construction will ever be needed in any scenario.

\subsubsection{From Single-Node Efficiency to Network-Level Efficiency}
\label{subsubsec:ch1-isac-core value-deployment-network}

The deployment advantage of ISAC is not limited to simplifying individual nodes; more importantly, it can generate significant scalability and evolution benefits at the network level.

\paragraph{1) Efficiency gain at the single-node level}
At the single-site level, an ISAC node can simultaneously provide communication and partial sensing functions, thereby reducing the number of devices, installation procedures, and duplicated configurations in antennas, power supply, cabinet space, and backhaul access. In scenarios with limited pole or tower space and dense roadside infrastructure, such functional integration can effectively reduce deployment and maintenance burden.

\paragraph{2) Native support for multi-site collaboration}
At the network level, the existing mechanisms in communication systems, such as time synchronization, frequency synchronization, inter-site interfaces, and resource coordination, can also support multi-site cooperative sensing \cite{Zhang2021ISAC6G,Liu2022ISAC}. Therefore, compared with building an independent cooperative sensing network from scratch, ISAC can reuse the existing infrastructure and significantly reduce the complexity of network-level deployment.

\paragraph{3) Easier scaling from local trial to wide-area deployment}
ISAC also enables a smooth evolution path from local trials to wide-area deployment. Operators may first activate sensing functions in hotspot areas or selected application scenarios, and then gradually expand them after validating their technical feasibility and service value. This incremental approach supports phased investment, deployment-time optimization, and reuse of the existing operation and maintenance framework, thereby reducing upgrade risks.

\paragraph{4) From engineering efficiency to ecosystem efficiency}
More broadly, the value of ISAC lies not only in engineering efficiency but also in its potential for platform-based expansion. With a unified network platform, equipment vendors, operators, and vertical industry users can collaborate around common interfaces, common management frameworks, and shared data flows \cite{IMT2030Framework,Letaief2019Roadmap}. This helps promote a more standardized and scalable deployment pathway.

\paragraph{5) A necessary caveat}
It should be noted that ISAC does not imply that all sensing tasks can be directly supported by the existing communication network. Its sensing capability still depends on factors such as network topology, operating frequency, bandwidth, site layout, and environmental blockage. In scenarios requiring high-precision imaging or strong robustness, ISAC may still need to operate jointly with cameras, radar, LiDAR, or other sensors \cite{Zhang2021ISAC6G, Wei2023Integrated}.

In summary, the deployment efficiency of ISAC is reflected in multiple dimensions: reducing duplicated construction at the single-node level, reusing existing coordination mechanisms at the network level, supporting gradual upgrades along the evolution path, and enabling platform-based expansion at the ecosystem level. These advantages make ISAC a particularly promising approach for rapidly introducing sensing capability into future wireless networks.

\subsection{Advantage in Business Closed-Loop Capability}
\label{subsec:ch1-isac-core value-closed loop}

Unlike conventional networks that mainly serve as information delivery infrastructures, Integrated Sensing and Communications (ISAC) not only provides data transmission capability but also enables real-time perception of the external environment, target states, and service objects. In this way, ISAC can provide upper-layer applications with timely and direct state information, thereby supporting a complete business closed loop of sensing, communication, decision-making, execution, and feedback \cite{Liu2022ISAC, Zhang2021ISAC6G, ITU2030, IMT2030Framework}.

From this perspective, ISAC drives networks to evolve from mere connectivity tools into platforms that support closed-loop service operation. The network is no longer limited to carrying service data, but can participate more deeply in the service process itself. For 6G application scenarios such as intelligent transportation, low-altitude economy, industrial automation, smart campuses, and public safety, high-rate and low-latency communication alone is insufficient to support highly autonomous, real-time, and reliable operations. Networks are also expected to possess awareness of the physical world and to closely cooperate with computing and control functions \cite{Letaief2019Roadmap, Saad2020Vision6G, ITU2030}.

As service paradigms shift from the traditional human-driven mode of “human initiation, human decision, and human execution” toward a new mode characterized by “autonomous device operation, real-time system coordination, and network-assisted control,” business closed-loop capability is becoming a key feature of future networks evolving from connectivity provision to intelligent service enablement.

\subsubsection{The limitations of traditional communication systems}
\label{subsubsec:ch1-isac-core value-closed loop-communication}

The primary goal of traditional communication systems is to enable reliable information delivery among users, terminals, devices, and service platforms. Their design is typically centered on communication-oriented metrics such as bandwidth, latency, reliability, coverage, connection density, and spectral efficiency \cite{Goldsmith2005,Andrews2014, Shafi20175G}. From the perspective of communication engineering, this paradigm has become highly mature and has achieved great success in mobile Internet, Internet of Things (IoT), and broadband access applications. However, when examined from the viewpoint of \emph{business closed-loop operation} rather than pure \emph{information transmission}, traditional communication systems still show a fundamental limitation: \textbf{they are effective at delivering information, but they usually do not directly provide awareness of the external physical environment}.

More specifically, although traditional communication networks can carry data generated by cameras, radars, industrial sensors, vehicular terminals, and application platforms, the network itself usually does not directly sense environmental states. As a result, it is difficult for the network to actively identify target locations, motion trajectories, behavior changes, space occupancy, or event evolution. In other words, traditional communication systems are good at answering questions such as \emph{where information should be delivered}, \emph{how fast it should be transmitted}, and \emph{how reliably it should arrive}, but they are generally not designed to answer questions such as \emph{what is happening in the environment}, \emph{where the target is}, \emph{how the state is changing}, and \emph{what action should be taken next}.

From a system architecture perspective, environmental perception and state estimation in traditional networks are usually supported by external sensing devices and separate processing platforms. For example, in vehicle--road coordination, environmental perception often relies on cameras, millimeter-wave radars, and LiDARs. In industrial automation, equipment monitoring and workstation identification are commonly supported by dedicated sensor networks and control buses. In low-altitude economy scenarios, target detection, trajectory identification, and regional warning often depend on independent surveillance and detection infrastructures. In such systems, the communication network mainly serves as a channel for data uploading, command delivery, and link maintenance, while environmental understanding, state estimation, and decision making are carried out outside the network.

This separated architecture, which can be summarized as \emph{external sensing, intermediate communication, and subsequent decision making}, is acceptable when system scale is limited and functions are relatively simple. However, in future scenarios requiring high real-time performance, strong autonomy, and large-scale coordination, it exposes several structural limitations.

\paragraph{1) Long processing chain and additional latency}
In the traditional mode, a complete service procedure usually involves the following chain: data acquisition by external sensors, data transmission through the communication network, analysis at the edge or cloud, generation of control commands, and command delivery back to the executor. Since sensing, transmission, computation, decision making, and execution are distributed across multiple heterogeneous subsystems, the end-to-end processing chain becomes long and may introduce extra transmission delay, processing delay, and interaction overhead \cite{Letaief2019Roadmap, Saad2020Vision6G}. In highly time-sensitive scenarios such as autonomous driving, collaborative robotics, and emergency response, this long chain can significantly reduce closed-loop responsiveness.

\paragraph{2) Complex interfaces and weak cross-system coordination}
In traditional architectures, communication, sensing, and control systems are often implemented by different devices, protocols, and management platforms. This leads to heterogeneous interfaces, inconsistent data formats, and complicated synchronization mechanisms, which increase the difficulty of system integration and cross-domain coordination. Even if each individual subsystem performs well, the overall business closed loop may still be constrained by interface bottlenecks, synchronization errors, and information loss.

\paragraph{3) Lack of global situational awareness at the network side}
Because traditional communication networks do not directly participate in environmental sensing, the network layer usually has access only to traffic status, connectivity status, and resource utilization, while lacking direct knowledge of target states, scene conditions, and task evolution. As a result, network optimization is typically based on communication-side objectives such as throughput, latency, and load balancing, rather than on task-oriented or situation-aware objectives. For example, the network may know that a roadside unit is heavily loaded, but it may not know whether there is an abnormal pedestrian crossing, a sudden vehicle stop, or a low-altitude intruder nearby. Therefore, it is difficult for the network to provide truly task-aware resource allocation and service support.

\paragraph{4) Difficult to support autonomous and closed-loop control}
Many future applications require the network to participate directly in closed-loop control rather than acting only as a command delivery channel. Examples include motion control in industrial Internet, cooperative collision avoidance in intelligent transportation, and trajectory conflict management in low-altitude networks \cite{Saad2020Vision6G, Liu2022ISAC}. These applications require the capability of \emph{sense, analyze, decide, and respond} in a tightly coupled manner. If the network itself has no sensing capability, the closed loop must rely on external systems for environmental information, which reduces response speed, stability, and autonomy. In this sense, traditional communication systems are naturally more suitable for open-loop information services than for strongly real-time autonomous closed-loop tasks.

\paragraph{5) Limited capability for continuous optimization}
A business closed loop does not only require one-time sensing and one-time control; it also requires the system to continuously adjust resource allocation and decision strategies according to execution outcomes. Since traditional communication networks do not directly observe the physical world, they usually cannot independently perform continuous feedback optimization across the chain of \emph{sensing results -- task outcomes -- network scheduling strategies} \cite{IMT2030Framework, Letaief2019Roadmap}. Consequently, in complex applications, the network often plays a passive role in responding to upper-layer requests rather than actively improving end-to-end task performance.

For teaching purposes, the limitations of traditional communication systems in supporting business closed loops can be summarized in Table~\ref{tab:traditional-network-closed-loop-limitations}.
\begin{table}[htbp]
\centering
\caption{Limitations of Traditional Communication Systems in Supporting Business Closed Loops}
\label{tab:traditional-network-closed-loop-limitations}
\begin{tabular}{p{3.6cm}p{4.8cm}p{4.8cm}}
\hline
\textbf{Aspect} & \textbf{Traditional Communication Network} & \textbf{Impact on Business Closed Loop} \\
\hline
Environment awareness & Usually relies on external sensors & Network lacks direct perception of physical states \\
System architecture & Sensing, communication and control are separated & Long processing chain and slow feedback \\
Coordination mechanism & Multi-system interfaces are heterogeneous & High integration complexity and weak coordination \\
Network intelligence basis & Mainly optimizes traffic and connectivity metrics & Difficult to optimize for task-oriented objectives \\
Control support capability & Mostly serves as a data transport channel & Limited support for autonomous real-time control \\
\hline
\end{tabular}
\end{table}

In summary, the fundamental limitation of traditional communication systems is not insufficient communication capability, but rather that their system role is primarily defined as an \textbf{information pipeline}, instead of being an integral part of \textbf{environmental awareness and business closed-loop operation}. As future networks evolve toward the integration of sensing, communication, computation, and control, this traditional positioning becomes increasingly inadequate for applications requiring high real-time performance, strong autonomy, and deep coordination \cite{Liu2022ISAC, ITU2030, Saad2020Vision6G}. This is exactly why the value of ISAC lies not only in improving spectral efficiency or reducing hardware redundancy, but also in enabling the network to combine communication capability and environmental awareness within a unified framework, thereby laying the foundation for truly closed-loop intelligent networking.

\subsubsection{ISAC Enables a Closed Loop of ``Perception--Transmission--Decision--Feedback''}
\label{subsubsec:ch1-isac-core value-closed loop-isac}

One of the key system-level values of ISAC is that it enables a business closed loop of \emph{perception--transmission--decision--execution--feedback}, allowing the network to evolve from a pure information delivery platform into an intelligent infrastructure capable of environmental awareness, task support, and adaptive optimization \cite{Liu2022ISAC, Zhang2021ISAC6G, Saad2020Vision6G, ITU2030}.  
Unlike the traditional architecture, where sensing, communications, decision-making, and control are largely separated, ISAC embeds sensing capability into the communication network itself. As a result, the chain from state acquisition to information sharing, decision-making, action execution, and performance feedback is significantly shortened, improving system timeliness, autonomy, and robustness.

From a system operation perspective, this closed loop typically consists of five stages: \textbf{perception}, \textbf{transmission}, \textbf{decision}, \textbf{execution}, and \textbf{feedback}.

\begin{enumerate}[label=(\arabic*)]
\item \textbf{Perception: acquiring target and environment states}

ISAC exploits integrated signal design, shared spectrum, and common infrastructure to sense target and environmental information during the communication process, such as distance, velocity, angle, position, trajectory, and scene variation \cite{Sturm2011WaveformDesign, Hassanien2016signaling, Liu2022ISAC}.  
Its key idea is that the network not only knows \emph{who is communicating}, but also perceives \emph{what is happening in the physical environment}.

\item \textbf{Transmission: sharing sensed states and delivering control information}

The sensed information must be promptly delivered to relevant nodes, edge platforms, or control centers, while the generated control commands must be reliably transmitted to the execution entities. Therefore, in the closed loop, communication supports three functions simultaneously: state synchronization, multi-node coordination, and command delivery. In this sense, the communication link serves as the information backbone of the entire closed-loop process \cite{ITU2030, IMT2030Framework}.

\item \textbf{Decision: performing scheduling, planning, and control}

Based on the acquired environmental states, the network side or application side can further perform resource scheduling, risk assessment, path planning, and control optimization \cite{Letaief2019Roadmap, Saad2020Vision6G}. In ISAC, decision-making may take place not only at the application layer but also at the network layer; for example, beam patterns, resource allocation, and sensing strategies can be adaptively adjusted according to target locations and environmental dynamics.

\item \textbf{Execution: translating decisions into actions}

The generated decisions must then be translated into physical actions or network actions, such as vehicle braking, UAV trajectory adjustment, robot collision avoidance, beam reconfiguration, handover, or resource rescheduling \cite{Saad2020Vision6G}. This stage ensures that the closed loop is not limited to awareness and planning, but is ultimately reflected in actual system behavior.

\item \textbf{Feedback: observing outcomes and refining future actions}

After execution, ISAC can continuously sense the updated target and environmental states, and feed them back to the system for performance evaluation, error correction, and subsequent decision refinement \cite{Zhang2021ISAC6G}. This forms a dynamic loop:
\[
\text{Perception} \rightarrow \text{Transmission} \rightarrow \text{Decision} \rightarrow \text{Execution} \rightarrow \text{Feedback}.
\]
Compared with open-loop operation, such a mechanism is more effective in coping with environmental uncertainty, target mobility, and external disturbances, thereby enhancing system stability and adaptivity \cite{Saad2020Vision6G}.
\end{enumerate}

In summary, the closed-loop capability enabled by ISAC is not a simple superposition of sensing and communication functions. Instead, it arises from their deep integration, which connects state acquisition, information sharing, intelligent decision-making, action execution, and feedback refinement into a unified process \cite{Liu2022ISAC, Feng2020JointRadarCommunication}. For highly dynamic and mission-critical scenarios such as intelligent transportation, low-altitude airspace management, industrial automation, and public safety, this closed-loop support capability is one of the most distinctive system-level advantages of ISAC over conventional communication networks.

\subsubsection{The Significance of Closed-Loop Service Capability for Typical Applications}
\label{subsubsec:ch1-isac-core value-closed loop-applications}

The significance of ISAC lies not only in adding sensing capability to communication networks, but more importantly in enabling a closed service loop of \emph{sensing, communication, decision-making, action, and feedback}. In this way, the network evolves from a pure information delivery platform into an intelligent infrastructure deeply embedded in task execution and system operation \cite{Liu2022ISAC, Zhang2021ISAC6G, Saad2020Vision6G, ITU2030}.  
In many future application scenarios, such as intelligent transportation, the low-altitude economy, industrial automation, and public safety, service processes are not one-time information exchanges. Instead, they are continuous and dynamically adjusted closed-loop processes. Therefore, the effectiveness of these systems often depends more on closed-loop capability than on communication capacity alone.

From a system perspective, the value of closed-loop service capability can be summarized in three aspects.  
\textbf{First, it improves real-time responsiveness.} Since sensing, information delivery, and control can be coordinated within a unified network, the delay from state change to control response can be significantly reduced.  
\textbf{Second, it enhances autonomous operation.} When the network can directly acquire environmental states and support decision-making, many tasks no longer rely entirely on human intervention or on separate external systems.  
\textbf{Third, it reduces cross-system coordination cost.} By integrating communication, sensing, and control on a more unified platform, ISAC helps alleviate interface fragmentation, information silos, and the overhead of multi-system coordination.

This significance is particularly evident in several representative applications.  
\textbf{In vehicle--road cooperation scenarios}, ISAC allows the network not only to exchange information among vehicles and infrastructure, but also to directly support target detection, trajectory estimation, and traffic situation updating. This enables a closed loop for cooperative driving, risk warning, and traffic management, thereby improving both driving safety and road efficiency \cite{Kumari2018IEEE80211adRadar,Ma2026integrated,Feng2020JointRadarCommunication}.  
\textbf{In low-altitude economy scenarios}, ISAC can organically connect target detection, state reporting, trajectory prediction, conflict avoidance, and control signaling, thus supporting an integrated closed loop of airspace sensing, scheduling, and control. This is essential for improving the safety and efficiency of large-scale low-altitude operations \cite{Jiang2025integrated,Meng2023uav}.  
\textbf{In industrial automation scenarios}, ISAC enables a unified wireless infrastructure to perform both device-state sensing and industrial data transmission, thereby supporting a closed loop of production monitoring, control execution, and feedback optimization. This is particularly valuable for flexible manufacturing, mobile robot coordination, and digitalized workshops \cite{Akyildiz2020Industrial6G,Kaushik2024integrated,Sturm2011WaveformDesign}.  
\textbf{In public safety scenarios}, ISAC contributes to a rapid closed loop of anomaly detection, information aggregation, collaborative decision-making, and response execution, thus improving situational awareness, risk warning, and coordinated emergency handling \cite{Ma2022WiFiSensingSurvey,Liu2022ISAC,Zhang2021ISAC6G}.  
\textbf{In smart home and intelligent space scenarios}, ISAC allows the network to continuously adjust services according to user behavior and environmental states, promoting wireless systems from merely \emph{connecting devices} to \emph{understanding users and environments} in a service-oriented closed loop \cite{Ma2022WiFiSensingSurvey,Zhang2020device}.

Overall, the fundamental importance of closed-loop service capability is that it allows the network to be more deeply embedded into the service process itself. This drives the evolution of the system from \emph{perceiving the environment} to \emph{understanding the environment}, from \emph{delivering information} to \emph{supporting decisions}, and further to \emph{executing actions} and \emph{continuously optimizing outcomes}. From a pedagogical and methodological perspective, this is also a key distinction between ISAC and the simple juxtaposition of separate communication and sensing systems: the goal of ISAC is not merely functional coexistence, but the construction of a unified closed-loop intelligent infrastructure through joint signal, resource, and system design.

\subsubsection{The business closed-loop capability is the key for ISAC to achieve its 6G native capabilities}
\label{subsubsec:ch1-isac-core value-closed loop-6g}

From the perspective of network evolution, \textbf{business closed-loop capability} is the key that enables ISAC to evolve from mere \emph{functional integration} to an \emph{inherent capability} of 6G. Only when sensing results can continuously participate in service processes and contribute to network operation, resource scheduling, task execution, and feedback optimization can the integration of communication and sensing demonstrate true system-level value \cite{Liu2022ISAC,Zhang2021ISAC6G,Saad2020Vision6G}.

First, \textbf{business closed-loop capability transforms sensing from an auxiliary measurement function into an essential network service capability}. In traditional communication systems, the primary objective is information delivery, while ranging, localization, target detection, and environmental observation are usually treated as additional functions \cite{Sturm2011WaveformDesign}. In contrast, under the ISAC framework, sensing results can directly support link establishment, beam management, resource allocation, mobility management, and application control \cite{Kaushik2024toward}. Therefore, sensing is no longer an optional add-on, but a fundamental capability for enabling highly reliable, low-latency, and environment-adaptive services. For 6G, this implies that the network must not only transmit bits, but also understand contexts, perceive states, and support actions \cite{Saad2020Vision6G,Strinati20216G}.

Second, \textbf{business closed-loop capability provides the foundation for ISAC to support native intelligence and system autonomy}. The evolution toward 6G is expected to move networks beyond connectivity enhancement toward intelligent operation, with greater abilities to understand the physical environment, user behavior, and service tasks, thereby enabling self-adaptation, self-optimization, and self-management \cite{Saad2020Vision6G,Strinati20216G}. In this process, the value of ISAC lies not only in spectrum sharing or hardware reuse, but more importantly in its ability to bring physical-world awareness into the network control loop. For example, the network may adjust beam directions and access strategies according to target location and motion, or optimize link configuration and handover decisions according to environmental changes \cite{Liu2022ISAC,Zhang2021ISAC6G}. Without such a closed-loop mechanism, sensing information would remain at the level of auxiliary observation and could hardly serve as an effective input for autonomous network control.

Third, \textbf{business closed-loop capability provides the interface basis for ``ISAC+'' and the further evolution toward integrated communication, sensing, computation, and control}. A major direction of 6G is the deep coordination of communication, sensing, computation, and control to support complex task-oriented services such as autonomous driving, low-altitude collaboration, and digital twins. If ISAC remains limited to waveform reuse, spectrum sharing, and hardware integration, its contribution is still mainly local optimization \cite{Sturm2011WaveformDesign}. Once equipped with business closed-loop capability, however, sensing outputs can naturally trigger edge computing, task understanding, and control execution, while execution outcomes can in turn be verified and refined through sensing feedback \cite{Shi2016EdgeComputing,Chen2024integrated}. In this sense, business closed-loop capability not only deepens the application value of ISAC itself, but also lays the methodological foundation for integrated communication-sensing-computation-control systems.

Finally, \textbf{business closed-loop capability determines whether ISAC can truly become a fundamental and core service capability in 6G networks}. The notions of \emph{native intelligence}, \emph{native sensing}, and \emph{inherent capability} in 6G do not simply refer to adding new functional modules; rather, they require these capabilities to be embedded into the network operation mechanism itself and to support service execution and system optimization \cite{Saad2020Vision6G,Strinati20216G}. For ISAC, even excellent ranging, detection, or localization performance is insufficient if sensing remains an isolated function. Only when ISAC forms a continuous cycle of sensing, transmission, decision-making, and feedback can it evolve from a technical concept into a native network capability \cite{Liu2022ISAC,Kaushik2024toward}.

Therefore, \textbf{business closed-loop capability not only reflects the essential advantage of ISAC over traditional communication systems and separated sensing systems, but also determines whether ISAC can be deeply integrated into the 6G architecture as a foundational enabler for future task-oriented services}. This also suggests that future ISAC research should not focus solely on parallel improvements in communication and sensing performance, but should further address how sensing information enters the network decision chain, participates in service control, and continuously creates system-level value through closed-loop operation.